\documentclass{anvil-uspto}

\title{Passive Thermal Compensation Network for a Piezoresistive Pressure Sensor}
\author{Dana R. Okoye \and Priya N. Venkatesan}

\begin{document}
\maketitle

%% ---------------------------------------------------------------------------
%% FIELD OF THE INVENTION
%% ---------------------------------------------------------------------------
\fieldoftheinvention

\anvilpara{The present disclosure relates to pressure sensors, and more
particularly to passive thermal compensation networks for piezoresistive
micro\-electro\-mechanical (MEMS) pressure sensors.}

%% ---------------------------------------------------------------------------
%% BACKGROUND OF THE INVENTION
%% ---------------------------------------------------------------------------
\background

\anvilpara{Piezoresistive MEMS pressure sensors convert mechanical strain in a
sensing diaphragm into a resistance change in a set of implanted piezoresistors,
typically arranged as a four-element Wheatstone bridge. Such sensors are
sensitive to temperature in two distinct ways. First, the temperature
coefficient of resistance (TCR) of the piezoresistors shifts the zero-pressure
operating point of the bridge, producing a zero-offset error. Second, the
temperature coefficient of sensitivity (TCS) of the diaphragm causes the bridge
span --- the output change per unit pressure --- to fall as temperature rises.}

\anvilpara{One approach measures temperature with a separate sensor and corrects
the pressure reading numerically in firmware using a stored calibration table.
Another approach places a resistor having a chosen temperature coefficient in
series with the bridge excitation supply to flatten the span-versus-temperature
curve. A third approach drives the bridge from a constant-current source to
reduce span drift to first order.}

\anvilpara{Each of the foregoing approaches leaves at least one of the two error
mechanisms uncorrected, or imposes a per-unit calibration over multiple
temperatures during production test. A passive compensation network that
corrects both span drift and offset drift, and that requires only a single
room-temperature trim, is therefore desirable.}

%% ---------------------------------------------------------------------------
%% SUMMARY OF THE INVENTION
%% ---------------------------------------------------------------------------
\summary

\anvilpara{Disclosed is a passive thermal compensation network for a
piezoresistive pressure sensor. A split-path excitation network divides the
bridge supply into a constant-current leg and a proportional-to-absolute-%
temperature (PTAT) leg, the relative contribution of the two legs being set by
a single ratio resistor \refnum{12}. The PTAT leg increases excitation with
temperature so as to cancel the negative temperature coefficient of sensitivity
of the bridge, while the constant-current leg fixes the low-temperature
operating point.}

\anvilpara{A self-referencing offset-cancellation node sums into the bridge
output a compensation signal taken from the midpoint of a matched dummy
half-bridge \refnum{20} fabricated adjacent to the sense bridge \refnum{10} and
exposed to the same die temperature and process corner. Because the dummy
half-bridge tracks the common-mode offset drift of the sense bridge, the
offset drift is cancelled without a per-unit trim.}

\anvilpara{A single span trim resistor \refnum{30} sets the overall gain at room
temperature and is trimmed exactly once, at approximately 25\,$^{\circ}$C,
during a single-point production test. Because the split-path excitation network
has already flattened the span-versus-temperature characteristic, the
room-temperature trim holds across the full operating range.}

%% ---------------------------------------------------------------------------
%% BRIEF DESCRIPTION OF THE DRAWINGS
%% ---------------------------------------------------------------------------
\briefdescriptionofdrawings

\anvilpara{FIG. 1 is a block diagram of a piezoresistive pressure sensor having
a passive thermal compensation network according to one embodiment.}

\anvilpara{FIG. 2 is a schematic diagram of the split-path excitation network of
FIG. 1.}

\anvilpara{FIG. 3 is a plan view of a MEMS die showing the placement of the
dummy half-bridge adjacent to the sense bridge.}

%% ---------------------------------------------------------------------------
%% DETAILED DESCRIPTION OF EMBODIMENTS
%% ---------------------------------------------------------------------------
\detaileddescription

\anvilpara{FIG. 1 shows a piezoresistive pressure sensor including a sense
bridge \refnum{10} formed of four implanted piezoresistors arranged as a
Wheatstone bridge on a sensing diaphragm, a split-path excitation network
\refnum{40}, a dummy half-bridge \refnum{20}, a summing node \refnum{50}, and a
span trim resistor \refnum{30}. The split-path excitation network \refnum{40}
supplies the bridge; the summing node \refnum{50} combines the bridge output
with a compensation signal from the dummy half-bridge \refnum{20}; and the span
trim resistor \refnum{30} sets the output gain.}

\anvilpara{\emph{Split-path excitation network.} Referring to FIG. 2, the
split-path excitation network \refnum{40} comprises a constant-current source
\refnum{42} and a PTAT current source \refnum{44} whose outputs are summed into
the high node of the sense bridge \refnum{10}. The constant-current source
\refnum{42} may be implemented as a bandgap-referenced current mirror delivering
a fixed current independent of temperature. The PTAT current source \refnum{44}
may be implemented as a current mirror whose reference is the difference in
base-emitter voltage between two bipolar junctions operated at unequal current
densities, producing a current proportional to absolute temperature. The ratio
of the PTAT current to the constant current is set by a single ratio resistor
\refnum{12} placed in the PTAT reference leg: increasing the resistance of
\refnum{12} reduces the PTAT slope. To build this network, a person of ordinary
skill selects the constant-current magnitude to set the desired bridge
operating point at the cold end of the range, then selects \refnum{12} so that
the added PTAT current raises the excitation at the hot end of the range by the
fraction needed to offset the measured negative temperature coefficient of
sensitivity of the bridge.}

\anvilpara{In one embodiment the PTAT-to-constant-current ratio is tunable from
0.4 to 1.2 and is preferably 0.7 for a sense bridge whose temperature
coefficient of sensitivity is approximately $-0.18\,\%/^{\circ}\mathrm{C}$. The
bridge excitation current ranges from 0.1\,mA to 2\,mA and is preferably
0.5\,mA. The ratio resistor \refnum{12} preferably has a tolerance of 0.05\,\%
or better and is a thin-film resistor; tolerances looser than 0.1\,\% degrade
the cancellation and are not preferred. The same network narrows cleanly to a
$0\,^{\circ}\mathrm{C}$ to $+85\,^{\circ}\mathrm{C}$ commercial range, in which
case the ratio-resistor tolerance may be relaxed.}

\anvilpara{\emph{Self-referencing offset-cancellation node.} The dummy
half-bridge \refnum{20} comprises two piezoresistors fabricated from the same
implant as the sense bridge \refnum{10} but located on an unstrained region of
the die, so that the dummy resistors experience the same temperature and the
same process corner as the sense resistors but do not respond to pressure. The
midpoint of the dummy half-bridge \refnum{20} therefore produces a voltage that
tracks the common-mode, temperature-dependent offset of the sense bridge
\refnum{10} while remaining insensitive to pressure. This voltage is scaled by a
fixed divider \refnum{52} and summed into the bridge output at the summing node
\refnum{50}, subtracting the offset drift. Because the cancellation is
self-referencing --- both bridges drift together --- no per-unit offset trim and
no temperature soak are required.}

\anvilpara{\emph{Single trim-once span resistor.} The span trim resistor
\refnum{30} sets the overall output gain. It is trimmed exactly once, at
approximately 25\,$^{\circ}$C, during a single-point production test: the part is
held at room temperature, a known reference pressure is applied, and \refnum{30}
is adjusted until the output equals the target full-scale value. No second
temperature point is needed, because the split-path excitation network
\refnum{40} has already flattened the span-versus-temperature curve, so the gain
chosen at room temperature holds across the operating range. The span trim
resistor \refnum{30} is preferably a thin-film resistor with a tolerance of
0.05\,\% and a temperature coefficient of resistance of 25\,ppm/$^{\circ}$C or
less. As an alternative, the span trim may be realized as a laser-trimmed
resistor network on the sensor printed circuit board.}

\anvilpara{\emph{Discrete-resistor embodiment.} In one embodiment all
compensation elements --- the ratio resistor \refnum{12}, the fixed divider
\refnum{52}, and the span trim resistor \refnum{30} --- are discrete
surface-mount thin-film resistors mounted on the sensor printed circuit board,
and the dummy half-bridge \refnum{20} is a second, unpressurized MEMS die
co-packaged with the sense die in a common cavity so that both dies share the
package thermal environment.}

\anvilpara{\emph{Integrated embodiment.} In another embodiment, shown in plan
view in FIG. 3, the dummy half-bridge \refnum{20} and the compensation network
are integrated on the same MEMS die as the sense bridge \refnum{10}, using the
same piezoresistor implant step, so that the dummy resistors and the sense
resistors share a process corner exactly. The dummy half-bridge \refnum{20} is
placed immediately adjacent to the sense bridge \refnum{10}, outside the
flexing region of the diaphragm, to maximize thermal and process tracking while
remaining insensitive to applied pressure.}

\anvilpara{The foregoing embodiments are illustrative. The split-path excitation
network \refnum{40}, the self-referencing offset-cancellation node, and the
single trim-once span resistor \refnum{30} may be used together, as described,
or independently, and the disclosed ranges and alternatives may be combined in
any operative manner.}

\end{document}
